\section{Keterkaitan Buku Ajar dengan RPS Berbasis OBE}
\label{sec:keterkaitan-rps}

Buku ini bukan kumpulan materi yang disusun menurut selera penulis.
Seluruh isinya diturunkan dari satu dokumen: Rencana Pembelajaran
Semester (RPS) mata kuliah Bahasa Pemrograman Tingkat Rendah yang
disusun dengan pendekatan \eng{Outcome-Based Education}.

\subsection{Prinsip Outcome-Based Education}

\term{Outcome-Based Education} (OBE) adalah pendekatan penyusunan
pembelajaran yang berangkat dari \emph{hasil} yang harus dicapai,
bukan dari materi yang ingin disampaikan. Urutan perancangannya
terbalik dibandingkan pendekatan tradisional.

\begin{figure}[H]
    \centering
    \begin{tikzpicture}[
        node distance=0.55cm,
        alur/.style={
            draw=obeBlue, line width=0.8pt, fill=obeBlueLight,
            rounded corners=3pt, minimum width=4.4cm,
            minimum height=0.92cm, align=center, font=\small,
            inner sep=5pt
        },
        panah/.style={-{Stealth[length=2.5mm]}, line width=1pt, obeTeal}
    ]
        \node[alur] (a) {\textbf{1. CPL}\\apa yang harus dikuasai lulusan};
        \node[alur, below=of a] (b) {\textbf{2. CPMK}\\kompetensi pada mata kuliah};
        \node[alur, below=of b] (c) {\textbf{3. Sub-CPMK}\\indikator yang dapat diuji};
        \node[alur, below=of c, fill=obeAmberLight, draw=obeAmber]
            (d) {\textbf{4. Asesmen}\\instrumen pengukur indikator};
        \node[alur, below=of d, fill=obeGreenLight, draw=obeGreen]
            (e) {\textbf{5. Materi \& kegiatan}\\sarana mencapai indikator};

        \draw[panah] (a) -- (b);
        \draw[panah] (b) -- (c);
        \draw[panah] (c) -- (d);
        \draw[panah] (d) -- (e);
    \end{tikzpicture}
    \caption{Perancangan mundur (\eng{backward design}) pada OBE: dari
    capaian lulusan menuju asesmen, lalu materi.}
    \label{fig:backward-design}
\end{figure}

Perhatikan bahwa pada Gambar~\ref{fig:backward-design} materi berada
pada urutan terakhir. Konsekuensinya penting: sebuah materi hanya
dianggap perlu bila ia membantu mahasiswa mencapai indikator yang
sudah ditetapkan dan dapat diukur. Sebaliknya, setiap indikator harus
memiliki asesmen yang benar-benar mengukurnya.

\subsection{Jalur Capaian Pembelajaran}

Empat tingkat capaian yang dipakai pada mata kuliah ini adalah
sebagai berikut.

\begin{description}
    \item[CPL (Capaian Pembelajaran Lulusan)] Kompetensi yang harus
          dimiliki seluruh lulusan program studi. Mata kuliah ini
          dibebani dua CPL, yaitu \CPL{02} dan \CPL{03}.

    \item[CPMK (Capaian Pembelajaran Mata Kuliah)] Kompetensi
          spesifik yang harus dikuasai setelah menyelesaikan mata
          kuliah ini. Terdapat empat CPMK, yaitu \CPMK{01} sampai
          \CPMK{04}.

    \item[Sub-CPMK] Penjabaran CPMK menjadi indikator yang lebih
          terukur dan dapat diuji. Terdapat \dec{16} Sub-CPMK, dari
          \SubCPMK{1.1} sampai \SubCPMK{4.3}.

    \item[Materi dan aktivitas] Isi setiap bab, contoh program,
          praktikum, dan latihan yang menjadi sarana mencapai
          Sub-CPMK.
\end{description}

Hubungan keempat tingkat tersebut beserta asesmennya ditunjukkan pada
Tabel~\ref{tab:jalur-capaian}.

\begin{table}[H]
    \centering
    \caption{Jalur capaian pembelajaran dan asesmennya.}
    \label{tab:jalur-capaian}
    \small
    \begin{tabular}{@{}p{2.6cm}p{8.4cm}p{3.4cm}@{}}
        \toprule
        \textbf{Tingkat} & \textbf{Rumusan} & \textbf{Diukur oleh} \\
        \midrule
        \CPL{02} & Cara kerja sistem komputer dan perancangan algoritma
        & UTS, tugas \\
        \addlinespace
        \CPL{03} & Perancangan, analisis, dan penerapan algoritma pada
        bahasa tertentu & UAS, proyek \\
        \addlinespace
        \CPMK{01} & Arsitektur register dan manajemen memori &
        UTS, tugas, kuis \\
        \addlinespace
        \CPMK{02} & Hubungan instruksi mesin dengan perangkat keras &
        UTS, tugas, kuis \\
        \addlinespace
        \CPMK{03} & Menulis kode assembly aritmatika dan logika &
        UAS, praktikum, proyek \\
        \addlinespace
        \CPMK{04} & \eng{Debugging} level instruksi memori &
        UAS, praktikum, proyek \\
        \bottomrule
    \end{tabular}
\end{table}

\subsection{Keterkaitan dengan RPS Mata Kuliah}

RPS mata kuliah ini diselenggarakan dalam \dec{16} pertemuan:
\dec{14} pertemuan kuliah dan praktikum, ditambah Ujian Tengah
Semester pada minggu ke-8 dan Ujian Akhir Semester pada minggu ke-16.
Setiap pertemuan dipetakan ke Sub-CPMK tertentu. Buku ini mengikuti
pemetaan tersebut secara persis, sehingga mahasiswa dapat membaca
bab yang bersesuaian dengan minggu perkuliahan yang sedang berjalan.

\begin{table}[H]
    \centering
    \caption{Kesepadanan pertemuan, Sub-CPMK, dan bab buku.}
    \label{tab:kesepadanan-pertemuan}
    \small
    \begin{tabular}{@{}cp{5.2cm}p{1.7cm}p{4.0cm}@{}}
        \toprule
        \textbf{Minggu} & \textbf{Materi pokok} & \textbf{Sub-CPMK} &
        \textbf{Bab} \\
        \midrule
        1 & Pengantar assembly dan sistem bilangan & 1.1 & Bab~III \\
        \addlinespace
        2 & Arsitektur 8086: BIU dan EU & 1.2 & Bab~IV \\
        \addlinespace
        3 & Register dan \eng{flag} & 1.3 & Bab~V \\
        \addlinespace
        4 & Segmentasi memori dan alamat fisik & 1.4 & Bab~VI \\
        \addlinespace
        5 & Memulai assembly: TASM, TLINK, COM, EXE, direktif & 2.1 &
        Bab~VII \\
        \addlinespace
        6 & Instruksi transfer data & 2.2 & Bab~VIII \\
        \addlinespace
        7 & Interupsi BIOS/DOS & 2.3 & Bab~IX \\
        \addlinespace
        8 & \textbf{Ujian Tengah Semester} & 1.1--2.3 & Bab~III--IX \\
        \addlinespace
        9 & Operasi aritmatika & 3.1 & Bab~X \\
        \addlinespace
        10 & Operasi logika dan \eng{shift} & 3.2 & Bab~XI \\
        \addlinespace
        11 & Tujuh mode pengalamatan & 3.3 & Bab~XII \\
        \addlinespace
        12 & Percabangan dan perulangan & 3.4 & Bab~XIII \\
        \addlinespace
        13 & Stack, prosedur, dan macro & 3.5 & Bab~XIV \\
        \addlinespace
        14 & Array, string, dan pencetakan bilangan & 3.6 & Bab~XV \\
        \addlinespace
        15 & \eng{Debugging} level instruksi & 4.1--4.3 & Bab~XVI \\
        \addlinespace
        16 & \textbf{Ujian Akhir Semester} & 3.1--4.3 & Bab~X--XVI \\
        \bottomrule
    \end{tabular}
\end{table}

\begin{infobox}
Bab~I dan Bab~II tidak terikat pada satu pertemuan tertentu. Bab~I
dibaca sebelum perkuliahan dimulai; Bab~II dibaca seiring pertemuan
pertama sebagai bekal konseptual. Bab~XVII dibaca menjelang Ujian
Akhir Semester.
\end{infobox}
